\chapter{Introduction}
\label{ch:introduction}

\section{Context}
\label{sec:context}

Contemporary sports clubs have evolved into complex, data-intensive organisations whose
day-to-day operations bear little resemblance to the amateur structures of previous
decades. A modern professional club is no longer a single team supported by a coach and a
physiotherapist; it is a multi-discipline enterprise that simultaneously fields teams across
several sports, employs dozens of specialised staff, manages hundreds of contracted
athletes, and engages a global audience of supporters. Each of these activities continuously
generates operational data: training sessions and the load they impose, match fixtures and
the events that unfold within them, medical incidents and the rehabilitation journeys that
follow, scouting assessments of prospective signings, sponsorship negotiations, contract
renewals, and the constant stream of communication that binds these functions together.
A club of this kind is, in effect, a small enterprise whose product is sporting performance
and whose raw material is information: who is fit, who is in form, who the next opponent is,
what the medical staff have advised, and what the commercial department has agreed.

To appreciate why the management of this information has become a first-order engineering
problem, it is useful to situate the present work within the broader sports-technology
landscape. Over the past two decades the discipline now widely referred to as
\emph{sports analytics} has matured from a niche curiosity into a mainstream competitive
advantage. Professional clubs routinely instrument their athletes with wearable sensors that
sample heart rate, acceleration and positional data hundreds of times per second; they
record and tag every match with frame-accurate video; and they purchase event-level data
feeds that describe every pass, shot and tackle in a fixture. The volume and velocity of this
data have grown to the point where it can no longer be managed by hand, and the clubs that
extract value from it are precisely those that have invested in the software infrastructure to
collect, integrate and reason over it. Sport, in other words, has become a software-intensive
domain, and the quality of a club's information systems is now a genuine determinant of its
performance both on and off the field.

This transformation has been mirrored on the supporter side. Consumer-facing data
platforms such as SofaScore and ESPN have reset public expectations about how sporting
information should be presented: fixtures, results, league standings and minute-by-minute
event timelines are now expected to be available instantly, updated live as matches unfold,
and rendered through polished, responsive interfaces on any device. A supporter who can
watch a match update second-by-second on a free mobile application will not tolerate a club
portal that displays last week's results entered by hand. The bar for immediacy, breadth and
presentation has been set by the consumer web, and any platform that a club exposes to its
own staff and supporters is implicitly measured against it.

The working setting for this project is the multi-sport club model exemplified by
FC~Barcelona, an organisation that operates several professional teams under a single
governing body while sharing common medical, scouting, analytics and commercial
departments. This setting was chosen deliberately because it captures the full breadth of
the problem: the same governing entity must reason about distinct sports, each with its own
positions, formations, competition calendars and performance indicators, while still
enforcing a single coherent identity, a unified access-control policy and a shared analytical
backbone. A platform that can serve such an organisation must therefore be genuinely
multi-tenant across disciplines, role-aware down to the level of individual job functions, and
capable of ingesting and presenting real competition data rather than merely storing
hand-entered records. The multi-sport club is not an arbitrary stress test; it is the natural
habitat of the very fragmentation problem this project sets out to solve, because every
shared department---medical, scouting, analytics, commercial---must serve teams whose
sporting logic differs while still drawing on one common pool of people, contracts and money.

The technical context is equally important. The systems that successfully manage data at
this scale in other industries---banking, logistics, e-commerce---have converged on a
common architectural style: rather than a single large program, they are built as a federation
of small, independently deployable services, each owning a well-defined slice of the domain
and its data, communicating over network protocols, and coordinated by shared
infrastructure for discovery, configuration and routing. This \emph{cloud-native
microservices} style~\cite{microservices} is not a fashion but a response to a concrete set of
pressures that a multi-sport club exhibits in full: heterogeneous domains that evolve at
different rates, the need to isolate failures so that an outage in one area does not take down
the whole organisation, and the desirability of scaling read-heavy public surfaces
independently of write-heavy internal tooling. The live-match engine, for example, is a
bursty, real-time workload that behaves nothing like the steady, transactional load of the
medical records subsystem, and forcing both into a single deployable unit would mean that
neither could be tuned, scaled or released on its own terms.

In practice, the data that should knit these activities together is instead scattered across
disconnected spreadsheets, third-party SaaS subscriptions, e-mail threads and printed
forms. A player's injury record maintained by the medical department is invisible to the
coach planning the next training block; a scout's report on a prospect never reaches the
sport manager negotiating the transfer in a structured form; and a supporter following the
club has no window onto the same fixtures, results and squad information that the club
itself maintains internally. The cost of this fragmentation is not merely inconvenience: it is
duplicated effort, reconciliation errors, stale information driving real decisions, and the
practical impossibility of asking cross-domain questions---such as whether a particular
training regime correlates with a rise in soft-tissue injuries, or whether a scout's
recommendations have historically panned out. This fragmentation is the central motivation
for the \emph{Multi-Sport Club Management System} (MSCMS), delivered under the product
brand \emph{Sportify}: a single, cloud-native platform that consolidates the operational
concerns of a multi-sport club and exposes each of them through one consistent, secured
interface, while adopting the real-time, data-rich presentation philosophy that the
consumer web has made the baseline expectation.

\section{Motivation}
\label{sec:motivation}

The motivation behind MSCMS is to consolidate every operational concern of a multi-sport
club into one cohesive platform, in which each professional role is presented only with the
information relevant to its responsibilities, and in which every data point is traceable,
secured and analytically meaningful. The design draws directly on the real organisational
structure of a large European club: a single governing body managing multiple professional
teams across several sports, supported by shared medical, scouting, analytics and
sponsorship functions. Modelling a real club rather than an abstract one forces the system
to confront the genuine complexity of multi-discipline operation instead of solving an
artificially simplified problem. Concretely, the platform models a catalogue of fourteen
distinct professional roles---administrator, head coach, assistant coach, specific coach,
fitness coach, performance analyst, team doctor, physiotherapist, team manager, sport
manager, scout, sponsor, player and supporter---each of which sees a different slice of the
system. Designing for fourteen roles rather than the three or four that generic team tools
assume is itself a forcing function: it exposes exactly the access-control, navigation and
data-ownership questions that a real club must answer and that simpler products are free to
ignore.

A second, equally important motivation is the desire to move beyond static record-keeping
towards a living, real-time platform. Earlier generations of club-management software
treated data as something to be entered after the fact: a fixture was recorded once it had
been played, an injury was logged once treatment had begun. The expectations of modern
staff and supporters, shaped by data-rich consumer platforms such as SofaScore and ESPN,
are very different. Users now expect matches to update live, notifications to arrive the
moment an event occurs, and dashboards to reflect the current state of the club rather than
yesterday's. MSCMS is therefore motivated not only by the goal of unifying fragmented
data, but by the goal of doing so in a manner that is real-time, event-driven and
operationally trustworthy. This ambition manifests in three concrete capabilities that recur
throughout this report: a live-match engine that drives a fixture through its phases as it is
actually played, recording goals, substitutions and injuries against a phase-bounded clock;
an event-driven notification fabric in which a significant occurrence in one domain
automatically surfaces an alert in another, without the two domains being directly coupled;
and a real-time team chat that lets staff coordinate over a persistent connection. Each of
these is real-time not as a veneer but as a structural property of the architecture.

A third motivation is correctness and trustworthiness of data. For a system to serve as a
single source of truth, it is not enough to gather data in one place; the data must be
internally consistent and protected against both accidental corruption and unauthorised
disclosure. This motivates two design commitments that shape much of the system. The first
is that the backend, not the browser or any external feed, is the authoritative source for
every fact the platform presents: when real competition data is ingested from an external
provider it is normalised and stored internally, and the front-end reads only from the
backend, so that there is exactly one version of the truth and one place to audit it. The
second is that sensitive operations are guarded consistently: medical and financial data are
visible only to the roles entitled to see them, and a destructive or irreversible action---such
as changing a password---is verified against the identity provider before it is allowed to
proceed. Trust, in a system that a club would actually rely on, is not a feature bolted on at
the end; it is a property that must be designed in from the data model outward.

Finally, the project is motivated by sound engineering practice. A platform of this breadth,
spanning identity, squad management, competitions, live matches, medical workflows,
messaging, notifications and analytics, cannot reasonably be built or maintained as a single
monolithic application. The motivation to adopt a cloud-native microservices
architecture~\cite{microservices} arises directly from the need for independent
deployability, fault isolation between domains, and the ability to evolve one functional area
(for example the live-match engine) without destabilising another (for example the medical
records subsystem). Decomposing the system along domain boundaries also yields a clean
mapping between the organisational structure of the club and the structure of the
software---squad operations, competitions and matches, medical and fitness, analytics and
reporting, and communication each become their own service with their own database---which
makes the system easier to reason about, to test in isolation, and to assign to different
developers working in parallel. The choice of Spring Boot~\cite{springboot} for the backend,
the Spring Cloud~\cite{springcloud} ecosystem for the discovery, configuration and gateway
control plane, and a RESTful inter-service contract~\cite{restapi} complemented by
asynchronous events follows naturally from this motivation.

\section{Scope}
\label{sec:scope}

The scope of MSCMS spans the principal operational domains of a multi-sport club, each
realised as a coherent feature area within the platform:

\begin{enumerate}
  \item \textbf{Identity and Role Management} --- centralised authentication through a
        self-hosted identity provider~\cite{keycloak} issuing signed JSON Web
        Tokens~\cite{jwt}, with a rich catalogue of fourteen professional roles organised
        into permission tiers and enforced uniformly across the system at three layers: the
        identity provider, the API gateway, and the front-end route guards~\cite{rbac}.
  \item \textbf{Squad and Club Operations} --- player and staff rosters, player profiles
        with photographs and detailed attributes, contracts, transfers and call-ups, and a
        dedicated Club Hub presenting the club's identity, squad, trophy cabinet and recent
        form.
  \item \textbf{Competitions} --- database-backed competitions covering a domestic league
        (La~Liga, modelled as a double round-robin) and the newly formatted UEFA
        Champions League (a single thirty-six-team league phase generated algorithmically,
        followed by a two-legged knockout bracket), with standings, fixtures and results
        recomputed and presented per competition, and completed competitions automatically
        locked read-only.
  \item \textbf{Match Operations} --- fixture creation with a mandatory starting line-up
        locked at creation, a live-match engine that drives a match through its phases in
        real time and automatically goes live at the exact scheduled kickoff, in-match event
        recording bounded by the current phase, per-sport substitution limits, and read-only
        match details for completed fixtures.
  \item \textbf{Medical and Fitness} --- a structured injury-management workflow that
        carries each affected player through a staged recovery journey---injury, diagnosis,
        treatment, rehabilitation, recovery and fitness testing---driven by the injury status
        as a single source of truth, alongside general medical catalogues, organised within a
        dedicated Treatment Room view that separates currently-injured from recovered
        players.
  \item \textbf{Communication} --- real-time team group chat with group membership,
        invitations that must be accepted, and message reactions, delivered over a persistent
        WebSocket connection~\cite{websocket} with a polling fallback.
  \item \textbf{Notifications and Alerts} --- an event-driven notification fabric, fed by a
        message broker~\cite{kafka}, that surfaces role-aware, priority-tagged alerts through
        an in-application notification centre reachable from a top-bar bell.
  \item \textbf{Reports, Analytics and Machine Learning} --- aggregated dashboards,
        downloadable reports, scouting and sponsorship workflows, and decoupled
        machine-learning services for match-outcome prediction and player-rating estimation,
        consumed transparently through the interface.
\end{enumerate}

The scope explicitly includes the ingestion of real competition data~\cite{footballdata} so
that competitions, fixtures and results reflect genuine football rather than synthetic
placeholders, and it includes a fan-facing surface so that supporters can follow the club
through the same platform the club uses internally. The scope deliberately excludes building
a bespoke identity provider, a custom message broker or a proprietary data feed; instead the
system integrates mature, well-understood components for these concerns and concentrates
its engineering effort on the club-management domain itself. Equally, the scope is bounded to
the football discipline for competitions and live matches---the domain in which authoritative
external data is available---while the role model, squad operations and medical workflows are
designed to generalise across sports. The boundaries of the system are drawn so that the
project demonstrates depth in the domains it claims rather than shallow breadth across
domains it cannot fully support.

\section{Problem Definition}
\label{sec:problem}

A modern multi-sport club faces a constellation of operational problems that no single
off-the-shelf product addresses end-to-end:

\begin{itemize}
  \item \textbf{Fragmented data across departments.} Medical, coaching, scouting,
        analytics and commercial functions each maintain isolated spreadsheets, documents
        and SaaS subscriptions. A player's injury record is invisible to the coach planning
        the next session, and a scout's assessment never reaches the sport manager in a
        structured, actionable form. The same fact is entered repeatedly in different places,
        each copy drifts out of step with the others, and no copy is authoritative.
  \item \textbf{Multi-discipline complexity.} A single club may field several sports under
        one governing body, each with distinct positions, formations, substitution rules,
        competition formats and performance indicators that generic, single-sport team tools
        cannot model simultaneously. A product that hard-codes the assumptions of one
        sport cannot serve a club whose departments must reason about several at once.
  \item \textbf{Inconsistent role-based access.} Tools lacking fine-grained access
        control either over-share data---exposing salary or medical information through a
        public portal---or over-restrict it, preventing legitimate staff from performing
        their duties. Worse, where access control exists only in the user interface, it can be
        bypassed simply by calling the underlying interfaces directly, so the apparent
        protection is illusory.
  \item \textbf{Absence of real-time operation.} Conventional systems record events after
        the fact and offer no live view of an in-progress match, no immediate notification
        when a significant event occurs, and no real-time channel through which staff can
        coordinate. The information they present is always, by construction, out of date.
  \item \textbf{Unstructured medical workflows.} Injuries are typically tracked as free
        text with no enforced progression from diagnosis through treatment, rehabilitation
        and fitness testing back to availability, making it impossible to reason about a
        player's recovery state at a glance, or to aggregate recovery patterns across the
        squad.
  \item \textbf{Limited analytical insight.} Without a unified data layer, computing
        cross-domain indicators---such as win-rate against a given opponent, injury
        incidence per position, or the contribution of a particular scout---is impractical,
        because the data needed to correlate one domain with another never sits in the same
        place.
  \item \textbf{Weak commercial and transfer workflows.} Sponsorship offers, contract
        negotiations and transfer windows are usually managed over e-mail and attached
        documents, with no audit trail, no structured state, and no automatic notification
        of the relevant stakeholders when an offer changes hands.
  \item \textbf{Brittle deployments and inconsistent identity.} Monolithic legacy systems
        force every domain onto a single release cadence, so that a small change to one area
        risks the stability of all the others, and the absence of a central identity provider
        leaves each department maintaining its own accounts, slowing onboarding,
        off-boarding and audit, and multiplying the surface for credential mismanagement.
\end{itemize}

Taken together, these problems define the gap that MSCMS sets out to close: the lack of a
unified, real-time, role-aware and independently deployable platform that models the genuine
complexity of a multi-sport club while remaining trustworthy enough to serve as a single
source of truth. The problem is not that any one of these capabilities is unprecedented---live
scores, video analysis, team scheduling and identity management all exist as separate
products---but that no single system combines them, and the act of combining them is itself
the substantive engineering challenge, because it forces consistent identity, consistent
access control and consistent data semantics across domains that the market has hitherto
treated in isolation.

\section{Goals and Objectives}
\label{sec:goals}

In response to the problems defined above, MSCMS pursues the following concrete,
measurable objectives:

\begin{itemize}
  \item \textbf{Unify operations} into a single platform that spans identity, squad and
        club operations, competitions, live matches, medical and fitness management,
        communication, notifications and analytics, so that every role works against one
        consistent data set rather than a patchwork of disconnected tools.
  \item \textbf{Enforce role-based access control consistently} at multiple layers---the
        identity provider, the API gateway, and the front-end route guards---so that the same
        permission rules apply whether a user navigates through the interface or calls the
        underlying APIs directly~\cite{rbac}. The objective is concretely measurable: a
        request for a resource a role is not entitled to must be refused regardless of the
        path by which it arrives.
  \item \textbf{Decompose the backend} into independently deployable Spring
        Boot~\cite{springboot} microservices, discoverable through a service registry and
        reachable through a single API gateway, configured from a central configuration
        server~\cite{springcloud}, communicating over a RESTful contract~\cite{restapi} and
        asynchronous events~\cite{microservices}.
  \item \textbf{Model real multi-sport structure} by representing distinct sports, each
        with its own positions, formations, substitution limits and competition formats,
        taking a real club as the working example so that the model is validated against
        genuine rather than imagined complexity.
  \item \textbf{Ingest real competition data}~\cite{footballdata} so that competitions,
        fixtures, results and standings reflect genuine matches, with the backend acting as
        the single source of truth for all externally sourced data and the front-end never
        contacting the external provider directly.
  \item \textbf{Deliver real-time operation} through a live-match engine that drives a
        fixture through its phases as it is played, an event-driven notification fabric, and a
        real-time team chat delivered over a persistent connection with a graceful polling
        fallback.
  \item \textbf{Structure medical care} as an explicit, staged recovery journey, driven by
        the injury status as the single source of truth, that makes each player's availability
        and rehabilitation state immediately legible at the level of both the individual and
        the squad.
  \item \textbf{Centralise analytics} into role-aware dashboards and downloadable reports,
        augmented by machine-learning predictions consumed transparently through the
        interface without exposing the non-technical user to the mechanics of the models.
  \item \textbf{Guarantee independent deployability} so that each domain ships on its own
        release cadence behind backward-compatible API contracts, and so that the failure or
        redeployment of one service does not compromise the availability of the others.
  \item \textbf{Deliver a coherent, club-branded experience} under the Sportify identity,
        so that staff and supporters interact with a polished, unified product rather than an
        assemblage of disjoint internal tools.
\end{itemize}

\section{Solutions}
\label{sec:solutions}

\subsection{Traditional Solutions}
\label{subsec:traditional}

Historically, multi-sport clubs have coordinated their operations through a combination of
manual and semi-digital tools. Coaching staff maintain training plans and line-ups in
spreadsheets or printed binders; medical departments record injuries and treatments in
paper charts or local workbooks; and scouts compile player reports as documents that
circulate by e-mail. Sponsorship and transfer negotiations are conducted over e-mail
threads with attached contracts, while supporter communication is handled separately
through generic websites and social-media channels. Where dedicated software is adopted,
it is typically single-sport and single-purpose---a video-analysis suite here, a fixture
scheduler there---so that the club ends up operating a patchwork of unconnected products,
each with its own login, its own data model and its own export format.

This patchwork places a heavy coordination burden on administrators: data is duplicated
across systems, reconciliation is manual, and any change to a player's status---an injury,
a transfer, a contract renewal---requires parallel updates across disconnected silos. The
approach scales poorly with the number of disciplines and makes cross-domain analytics
effectively impossible, because no single store holds the data required to correlate, for
example, training load with injury frequency, or scouting activity with transfer outcomes.
Crucially, traditional approaches also lack consistent role-based access control: sensitive
medical or financial data is either over-shared through common drives or over-restricted to
the point that legitimate staff cannot perform their work. None of these arrangements offers
a live view of an in-progress match, an immediate notification when an event occurs, or a
structured medical workflow, and none presents a coherent, club-branded surface to
supporters. The cumulative effect is an organisation that spends a substantial fraction of its
administrative effort simply keeping copies of the same information in agreement with one
another---effort that produces no sporting value and that grows worse, not better, as the
club adds disciplines and staff.

\subsection{Technological Solution}
\label{subsec:technological}

To overcome these limitations, MSCMS (Sportify) was developed as a cloud-native,
multi-sport club-management platform built on a modern microservices
architecture~\cite{microservices}. The system unifies every operational concern of a
multi-discipline club into a single coherent interface, while preserving the independent
deployability of each functional domain. The backend comprises six Spring
Boot~\cite{springboot} domain microservices---user management, player management,
training-and-match, medical-and-fitness, reports-and-analytics, and notification-and-mail
---supported by a Spring Cloud~\cite{springcloud} control plane of a Eureka discovery
server, a centralised configuration server and an API gateway, and complemented by
dedicated, decoupled machine-learning services for match-outcome prediction and
player-rating estimation. Every domain service is registered with the discovery server, draws
its configuration from the central configuration server, and is reached through the single API
gateway, which performs token verification, request routing and cross-origin handling.
Services expose RESTful interfaces~\cite{restapi} for synchronous calls and emit
asynchronous events through a message broker~\cite{kafka} for cross-domain reactions, so
that, for instance, a goal recorded by the live-match engine automatically produces a
notification without any direct coupling between the two domains.

Authentication is delegated to a self-hosted identity provider~\cite{keycloak} that owns a
dedicated realm and manages the full catalogue of fourteen professional roles spanning club
leadership, coaching staff, medical and fitness staff, analytics and scouting roles, commercial
roles, and end-users such as players and supporters. Access tokens issued by the identity
provider take the form of signed JSON Web Tokens~\cite{jwt}; they are verified at the
gateway and decoded by the front-end so that the same permission map is enforced on both
layers~\cite{rbac}, and a request for an unauthorised resource is refused whether it arrives
through the user interface or through a direct API call. Each microservice owns its own
relational database~\cite{postgres}, following the database-per-service pattern, while
asynchronous events flow through the message broker and a shared cache accelerates
read-heavy paths. The front-end is a modern single-page application built with
Next.js~\cite{nextjs} and React~\cite{react} that presents role-aware navigation, a
fan-facing read-only surface, real-time match and chat views, and aggregated analytics
dashboards. The entire stack is containerised and orchestrated with Docker~\cite{docker},
so that the whole platform can be brought up reproducibly and self-hosted by the club.
The result is a scalable, observable and role-aware platform that directly addresses the
fragmentation, multi-discipline complexity, absence of real-time operation and analytical
blindness that characterise the traditional approaches described above, delivering them
within one cohesive, self-hostable system.
